\documentclass[12pt]{article}

%%%%%%%%%%%%%%%%%%%%%%%%%%%%%%%%%%%%%%%%%%%%%%%%%%%%%%%%%%%%%%%%%%
%                     PREAMBLE / PACKAGES
%%%%%%%%%%%%%%%%%%%%%%%%%%%%%%%%%%%%%%%%%%%%%%%%%%%%%%%%%%%%%%%%%%

% ----------------------------------------------------------------
% Font encoding and page geometry
% ----------------------------------------------------------------
\usepackage[T1]{fontenc}
\usepackage[margin=1in]{geometry}

% ----------------------------------------------------------------
% Line and paragraph spacing
% ----------------------------------------------------------------
\usepackage{setspace}
\usepackage{parskip}

% ----------------------------------------------------------------
% Section and table-of-contents formatting
% ----------------------------------------------------------------
\usepackage{titlesec}
\usepackage{tocloft}

% ----------------------------------------------------------------
% Graphics, figures, and page layouts
% ----------------------------------------------------------------
\usepackage{graphicx}
\usepackage{animate} % frame-by-frame animations in supported PDF viewers
\usepackage[font=small]{caption}
\usepackage{subcaption}
\usepackage{float}
\usepackage{placeins}
\usepackage{wrapfig}
\usepackage{multicol}

% ----------------------------------------------------------------
% Mathematics
% ----------------------------------------------------------------
\usepackage{amsmath}
\usepackage{amssymb}
\usepackage{amsthm}
\usepackage{mathrsfs}
\usepackage{cancel}
\usepackage{gensymb}

% ----------------------------------------------------------------
% Tables and lists
% ----------------------------------------------------------------
\usepackage{booktabs}
\usepackage{tabularx}
\usepackage{multirow}
\usepackage{enumitem}

% ----------------------------------------------------------------
% Colors and highlighting
% ----------------------------------------------------------------
\usepackage{xcolor}
\usepackage{soul}
\sethlcolor{yellow}

% ----------------------------------------------------------------
% Headers and footers
% ----------------------------------------------------------------
\usepackage{fancyhdr}

% ----------------------------------------------------------------
% Source-code examples
% ----------------------------------------------------------------
\usepackage{listings}

% ----------------------------------------------------------------
% Citations and hyperlinks
% ----------------------------------------------------------------
\usepackage[numbers]{natbib}
\usepackage{hyperref}

% ----------------------------------------------------------------
% Default line spacing
% ----------------------------------------------------------------
\onehalfspacing


\title{Lab 8}
\author{Chuan Wei Heng}
\date{\today}

\begin{document}

\maketitle

\section*{3.1 Constructing Piecewise Linear Approximations}

We begin by constructing a piecewise linear spline
interpolator. The idea is rather simple: given a set of
interpolation nodes, we construct a straight line between
each pair of adjacent nodes. The resulting piecewise
function is continuous, although its first derivative
is generally discontinuous at the interpolation nodes.

We first test the code on the function

\[
f(x)=e^x,\qquad x\in[0,1].
\]

The domain was divided into 10 equally spaced intervals,
and the resulting approximation is shown below.

\begin{figure}[H]
    \centering
    \includegraphics[width=0.85\linewidth]{figure/3.1.png}
    \caption{Linear spline interpolation of $f(x)=e^x$
    using 10 uniformly spaced intervals, together with
    the corresponding absolute error.}
    \label{fig:linear_exp}
\end{figure}

We can see that the linear spline approximates the
exponential function reasonably well. Interestingly,
the error vanishes at the interpolation nodes, which
makes sense since the spline is constructed to pass
through these points. Between nodes, the error increases
before returning to zero at the next interpolation point.


\section*{3.2 Linear Spline Experiment}

Next, we consider the function

\[
f(x)=\frac{1}{1+100x^2},\qquad x\in[-1,1].
\]

Unfortunately, Lab 7 was not even posted, so I am unable
to reproduce the experiments referenced in the assignment.
As such, I followed the procedure provided by Professor
Gillman in her construction of the linear spline code,
using 10 uniformly spaced intervals and comparing the
interpolant against the exact function.

\begin{figure}[H]
    \centering
    \includegraphics[width=0.85\linewidth]{figure/3.2.png}
    \caption{Linear spline approximation of the
    function using 10 uniformly spaced intervals, together
    with the corresponding absolute error.}
    \label{fig:linear_runge}
\end{figure}

We can see that the approximation is less accurate near
the center of the domain, where the function changes
rapidly. The maximum absolute error was approximately

\[
\|f-S_{\mathrm{linear}}\|_{\infty}
\approx 1.10\times10^{-1}.
\]

Nevertheless, the interpolant remains relatively well
behaved and does not exhibit the large oscillations
commonly associated with high-degree global polynomial
interpolation using uniformly spaced nodes.


\section*{3.3 Constructing Natural Cubic Splines}

We now extend the linear spline construction to cubic
splines. Rather than connecting neighboring interpolation
nodes using straight lines, we construct a cubic polynomial
on each interval.

Unlike linear splines, we additionally require continuity
of the first and second derivatives at the interpolation
nodes. This produces a smoother interpolant.

For this lab, we construct a natural cubic spline, which
satisfies the boundary conditions

\[
S''(a)=S''(b)=0.
\]

Following the formulation provided in the lab, the local
cubic polynomial is written as

\[
S_i(x)=
\frac{
M_i(x_{i+1}-x)^3+
M_{i+1}(x-x_i)^3
}{6h_i}
+C_i(x_{i+1}-x)+D_i(x-x_i),
\]

where $h_i=x_{i+1}-x_i$ and $M_i=S''(x_i)$.

The coefficients $M_i$ are obtained by constructing and
solving a tridiagonal linear system, which follows from
enforcing continuity of the first derivative. Once these
coefficients are known, the remaining coefficients
$C_i$ and $D_i$ are evaluated directly.

The resulting spline can then be evaluated on each
subinterval using the corresponding local cubic polynomial.


\section*{3.4 Cubic Spline Experiment}

Finally, we apply the natural cubic spline interpolator
to the same function considered in Section 3.2:

\[
f(x)=\frac{1}{1+100x^2},
\qquad x\in[-1,1].
\]

We again use 10 uniformly spaced intervals and compare
the resulting cubic spline against the linear spline
constructed previously.

\begin{figure}[H]
    \centering
    \includegraphics[width=0.85\linewidth]{figure/3.4.png}
    \caption{Comparison between linear and natural cubic
    spline interpolation of the function using
    10 uniformly spaced intervals.}
    \label{fig:cubic_runge}
\end{figure}

Interestingly, the linear spline actually performs better
than the natural cubic spline for this particular example.
While the cubic spline is smoother, it exhibits noticeable
oscillations between some interpolation nodes, particularly
near the center of the domain.

The maximum absolute errors were approximately

\[
\begin{aligned}
\|f-S_{\mathrm{linear}}\|_{\infty}
&\approx 1.10\times10^{-1},\\
\|f-S_{\mathrm{cubic}}\|_{\infty}
&\approx 1.90\times10^{-1}.
\end{aligned}
\]

This initially seems counterintuitive, since cubic splines
are generally expected to provide more accurate approximations
than linear splines. However, the function varies rapidly
near $x=0$, and the relatively coarse interpolation grid
does not adequately resolve this behavior.

As such, the additional smoothness constraints imposed on
the cubic spline do not necessarily translate into a
more accurate approximation. With a sufficiently refined
grid, we would generally expect the cubic spline to
perform substantially better.


\section*{Conclusion}

Overall, in this lab, we constructed both linear and
natural cubic spline interpolators. The linear spline
was relatively straightforward to implement, whereas
the cubic spline required solving a linear system to
enforce the necessary smoothness conditions.

Interestingly, we observed that the linear spline
outperformed the cubic spline when approximating the function 
$\frac{1}{1+(10x)^2}$ using only 10 intervals. This demonstrates
that a smoother interpolant is not necessarily a more
accurate one, particularly when the interpolation nodes
are too widely spaced.

For details of the code, we turn the reader towards the
Git repository by \textit{chuanweiheng-cell} entitled
APPM4600. Thank you!

\end{document}